\documentclass[a4paper, 12pt, twoside]{article}
\usepackage[utf8]{inputenc} % LaTeX, comprend les accents !
\usepackage[T1]{fontenc}
\usepackage[francais]{babel}
\usepackage{lmodern}
\usepackage{ae,aecompl}
\usepackage[top=2.5cm, bottom=2cm,
			left=3cm, right=2.5cm,
			headheight=15pt]{geometry}
\usepackage{graphicx}
\usepackage{eso-pic} % Nécessaire pour mettre des images en arrière plan
\usepackage{array}
\usepackage{hyperref}
\input{pagedegarde}


\title{Titre de mon projet}

\datedebut{jour mois anéée}
\datefin{jour mois année}


\membrea{Nom  prénom étudiant 1}
\membreb{Nom  prénom étudiant 2}



\begin{document}
\pagedegarde
\section*{Remerciements}
Je remercie mon chien, mon chat, mon cousin...
\newpage

\tableofcontents
\newpage

\section{Introduction}
\paragraph{Attention !} Vous devez impérativement utiliser ce template pour réaliser votre rapport de projet d'harmonisation. Vous n'avez pas le droit de changer le format (A4), la police (12pt), l'interligne ou les marges.\\

Voici un exemple pour citer un ouvrage ou un site web \cite{cat} et un exemple de tableau :

\begin{table}[h]
	\begin{center}
		\begin{tabular}{|c|c|}
			\hline
			• & • \\
			\hline
			• & • \\
			\hline
		\end{tabular}
	\end{center}
	\label{referencedutableau}

	\caption{Titre du tableau : légende du tableau}
\end{table}
Et ne jamais oublier que lorsqu'on place un tableau, on doit l'utiliser comme ici avec le tableau Table \ref{referencedutableau}. et on fait la meme chose avec une image, comme avec la Figure \ref{Tux}.

\begin{figure}[h]
	\centering
	\includegraphics{Tux.png}
	\caption{Tux, le pingouin}
	\label{Tux}
\end{figure}



\section{Environnement de travail}
\section{Description du projet et objectifs}
\subsection{Exemple de sous section}
\subsection{Un autre exemple de sous section}

\section{Bibliothèques, Outils et technologies}
\section{Travail réalisé}
\section{Difficultés rencontrées}
\section{Bilan}
\subsection{Conclusion}
\subsection{Perspectives}



\newpage
\section{Bibliographie}
\renewcommand{\bibname}{}
\renewcommand{\refname}{}
\begin{thebibliography}{2}
	\bibitem[label]{cle} Auteur, TITRE, editeur, annee
	\bibitem[LAM94]{lam1} L. LAMPORT, {\it \LaTeX : A Document preparation system, Addison-Wesley, 1994}
\end{thebibliography}

\newpage
\section{Webographie}
\begin{thebibliography}{2}
	\bibitem[CAT]{cat} \url{savoircoder.fr/cat}
\end{thebibliography}


\newpage
\section{Annexes}
\appendix
\makeatletter
\def\@seccntformat#1{Annexe~\csname the#1\endcsname:\quad}
\makeatother
\section{Cahier des charges}
\section{Exemple d'exécution du projet}
\section{Manuel utilisateur}


\end{document}